\chapter{الإشارات والأمواج}\label{ch:g10:signals-and-waves}

كل رسالة تلقّيتها يومًا وصلت إشارةً: صوتًا ضاغطًا
على طبلة أذنك، وهذه الصفحة بلغتك ومضةَ راديو. وخلف
التنوع فكرة واحدة --- \omterm{def:g10:orders-of-magnitude:unit}{مقدار فيزيائي} يتغير مع الزمن --- وعُدّة
واحدة: \omterm{def:g10:signals-and-waves:period}{الدور} \omterm{def:g10:signals-and-waves:frequency}{والتواتر} \omterm{def:g10:signals-and-waves:amplitude}{والسعة} وفنّ توقيت تأخير.
وبها يقيس قارب قاع البحر ويعدّ طبيب نبضات القلب.

\section{الإشارات تحمل المعلومات}

\begin{definition}[الإشارة]\label{def:g10:signals-and-waves:signal}
\emph{الإشارة}\index{إشارة} \omterm{def:g10:orders-of-magnitude:unit}{مقدار فيزيائي} يتغير مع الزمن
ليحمل معلومة: \omterm{def:g10:pressure:pressure}{كضغط} الهواء عند طبلة أذنك، أو التوتر عند
مربطي لاقط \omterm{def:g10:signals-and-waves:sound}{صوت}، أو سطوع الضوء في \omterm{def:g10:refraction:fiber}{ليف بصري}.
ولدراسة إشارة، سجّلها دالةً للزمن واقرأ
المعلومة من المنحني.
\end{definition}

في مكالمة هاتفية تغيّر المعلومة نفسها حاملها مرارًا:
ضغطًا في الهواء، ثم توترًا في اللاقط، ثم راديو عند
الهوائي --- وبالعكس عند الطرف الآخر. ومفردات واحدة تعالجها
كلها؛ ويبنيها هذا الفصل.

\section{الإشارات الدورية}

\begin{definition}[الإشارة الدورية، الدور]\label{def:g10:signals-and-waves:period}
تكون \omterm{def:g10:signals-and-waves:signal}{الإشارة} \emph{دورية}\index{إشارة دورية} إذا تكررت
تكرارًا متطابقًا على فترات منتظمة. و\emph{الدور}\index{دور} $T$ هو
مدة دورة كاملة واحدة، \omterm{def:g10:orders-of-magnitude:unit}{بالثانية}.
\end{definition}

\begin{definition}[التواتر]\label{def:g10:signals-and-waves:frequency}
\emph{تواتر}\index{تواتر} \omterm{def:g10:signals-and-waves:period}{إشارة دورية} هو عدد
الدورات في \omterm{def:g10:orders-of-magnitude:unit}{الثانية}:
\[
f = \frac{1}{T} .
\]
ووحدته، أي دورة واحدة في \omterm{def:g10:orders-of-magnitude:unit}{الثانية}، هي \emph{الهرتز}\index{هرتز}
(\unit{Hz})، مع المضاعفات المعتادة \unit{kHz} و\unit{MHz} و\unit{GHz}.
\end{definition}

\begin{definition}[السعة]\label{def:g10:signals-and-waves:amplitude}
\emph{سعة}\index{سعة} \omterm{def:g10:signals-and-waves:period}{إشارة دورية} هي أقصى انحراف
لها عن قيمة الراحة. ففي توتر يتذبذب تذبذبًا متناظرًا
حول الصفر، تكون السعة هي القيمة القصوى؛ أما من القمة إلى القاع فيقاس
\emph{ضعف} السعة (أي القيمة من قمة إلى قمة).
\end{definition}

\begin{example}[الشبكة ونغمة لا]\label{ex:g10:signals-and-waves:mains}
يتذبذب توتر الشبكة \omterm{def:g10:signals-and-waves:frequency}{بتواتر} $f = \qty{50}{Hz}$: فدوره
$T = 1/f = \qty{0.020}{s} = \qty{20}{ms}$. ونغمة لا القياسية في أوركسترا
\omterm{def:g10:signals-and-waves:signal}{إشارة} \omterm{def:g10:pressure:pressure}{ضغط} تواترها \qty{440}{Hz}: فتدوم الدورة الواحدة
$T = 1/440 = \qty{2.3}{ms}$. فكلما ارتفع \omterm{def:g10:signals-and-waves:frequency}{التواتر} قصرت الدورة: إذ إن $f$ و$T$
مقلوب أحدهما الآخر.
\end{example}

\begin{omfigure}
\begin{tikzpicture}[font=\small]
  \draw[-Stealth] (-0.3,0) -- (7.7,0) node[below] {$t$};
  \draw[-Stealth] (0,-1.7) -- (0,2.0) node[left] {إشارة};
  \draw[omDef, thick, smooth, samples=120, domain=0:7.3]
    plot (\x, {1.3*sin(120*\x)});
  \draw[dashed, thin] (0.75,1.3) -- (0.75,1.6) (3.75,1.3) -- (3.75,1.6);
  \draw[omThm, Stealth-Stealth] (0.75,1.6) -- (3.75,1.6)
    node[midway, above] {الدور $T$};
  \draw[omProp, Stealth-Stealth] (5.25,0) -- (5.25,-1.3)
    node[midway, right] {السعة};
\end{tikzpicture}

{\small \omterm{def:g10:signals-and-waves:period}{إشارة دورية}: يتكرر النمط كل $T$ \omterm{def:g10:orders-of-magnitude:unit}{ثانية}، \omterm{def:g10:signals-and-waves:amplitude}{والسعة}
هي أقصى انحراف عن قيمة الراحة.}
\end{omfigure}

\section{قراءة مخطط اهتزاز}

\begin{definition}[راسم الاهتزاز]\label{def:g10:signals-and-waves:oscilloscope}
\emph{راسم الاهتزاز}\index{راسم الاهتزاز} يرسم توترًا بدلالة الزمن؛
والمنحني على شاشته المشبَّكة هو \emph{مخطط اهتزاز}%
\index{مخطط اهتزاز}. وضبطان يحوّلان تقسيمات الشاشة إلى قيم
فيزيائية: \emph{قاعدة الزمن}\index{قاعدة الزمن} (بالثواني لكل تقسيمة
أفقية) و\emph{الكسب}\index{كسب} العمودي (بالفولت لكل تقسيمة
عمودية).
\end{definition}

\begin{method}[من الشاشة إلى الأعداد]\label{met:g10:signals-and-waves:oscillogram}
\begin{enumerate}
  \item عُدّ التقسيمات الأفقية التي تمتد عليها دورة كاملة واحدة
        (ومن قمة إلى قمة أسلم)؛ واضرب في \omterm{def:g10:signals-and-waves:oscilloscope}{قاعدة الزمن} للحصول على
        $T$، ثم $f = 1/T$.
  \item عُدّ التقسيمات العمودية من الخط الأوسط إلى قمة
        واضرب في \omterm{def:g10:signals-and-waves:oscilloscope}{الكسب} للحصول على \omterm{def:g10:signals-and-waves:amplitude}{السعة} (ومن قمة إلى قاع تحصل على
        القيمة من قمة إلى قمة: فنصّفها).
\end{enumerate}
\end{method}

\begin{example}[قراءة كاملة]\label{ex:g10:signals-and-waves:reading}
في \omterm{def:g10:signals-and-waves:oscilloscope}{مخطط الاهتزاز} أدناه، \omterm{def:g10:signals-and-waves:oscilloscope}{قاعدة الزمن} \qty{5}{ms} لكل تقسيمة
\omterm{def:g10:signals-and-waves:oscilloscope}{والكسب} \qty{2}{V} لكل تقسيمة. وتمتد الدورة الواحدة على $4.0$ تقسيمات:
$T = 4.0 \times \qty{5}{ms} = \qty{20}{ms}$ و
$f = 1/\qty{0.020}{s} = \qty{50}{Hz}$ --- أي الشبكة مرة أخرى. وتقع قمة على ارتفاع
$3.0$ تقسيمة فوق المحور: \omterm{def:g10:signals-and-waves:amplitude}{فالسعة}
$3.0 \times \qty{2}{V} = \qty{6.0}{V}$.
\end{example}

\begin{omfigure}
\begin{tikzpicture}[x=0.62cm, y=0.62cm, font=\small]
  \draw[gray!35, thin] (0,0) grid (10,8);
  \draw[gray!80] (0,4) -- (10,4);
  \draw[gray!80] (5,0) -- (5,8);
  \draw[omDef, thick, smooth, samples=120, domain=0:10]
    plot (\x, {4 + 3*sin(90*\x)});
  \draw[dashed, thin] (1,7) -- (1,7.6) (5,7) -- (5,7.6);
  \draw[omThm, Stealth-Stealth] (1,7.6) -- (5,7.6)
    node[midway, above] {$4.0$ تقسيمات};
  \draw[omProp, Stealth-Stealth] (7,4) -- (7,1)
    node[midway, right] {$3.0$ تقسيمات};
\end{tikzpicture}

{\small \omterm{def:g10:signals-and-waves:oscilloscope}{مخطط اهتزاز}، \omterm{def:g10:signals-and-waves:oscilloscope}{قاعدة الزمن} \qty{5}{ms}/تقسيمة \omterm{def:g10:signals-and-waves:oscilloscope}{والكسب} \qty{2}{V}/تقسيمة:
تمتد الدورة الواحدة على $4.0$ تقسيمات ($T = \qty{20}{ms}$، $f = \qty{50}{Hz}$)
وترتفع القمة $3.0$ تقسيمة (\omterm{def:g10:signals-and-waves:amplitude}{السعة} \qty{6.0}{V}).}
\end{omfigure}

\section{الصوت والإشارات الكهرمغناطيسية}

\begin{definition}[الصوت]\label{def:g10:signals-and-waves:sound}
\emph{الصوت}\index{صوت} \omterm{def:g10:signals-and-waves:signal}{إشارة} \omterm{def:g10:pressure:pressure}{ضغط}: فجسم مهتزّ ---
وترٌ أو غشاء أو حبال صوتية --- يدفع الهواء دفعًا إيقاعيًّا، وتنتقل
التضاغطات إلى الخارج بسرعة نحو \qty{340}{m/s}. ويحوّل لاقط الصوت
الضغطَ الواصل إلى توتر؛ وتحوّله طبلة الأذن إلى نبضات
عصبية. والصوت يحتاج وسطًا: فجرس يُقرع تحت جرس خلاء يصمت
كلما ضُخّ الهواء إلى الخارج.
\end{definition}

\begin{definition}[الموجة]\label{def:g10:signals-and-waves:wave}
الاضطراب المنتقل --- كنمط تضاغط \omterm{def:g10:signals-and-waves:sound}{الصوت}، أو التموّج
على بركة، أو نبضة راديو --- هو \emph{موجة}\index{موجة}. أما الوسط،
حين يوجد، فيبقى مكانه: إذ ترتفع كل رقعة ماء وتنخفض في مكانها بينما
يعبر التموّج البركة. أما هندسة الأمواج --- \omterm{def:g10:light-spectra:wavelength}{طول الموجة}
والتداخل والحيود --- فتُتناول في فصل لاحق؛ ولا يحتاج هذا
الفصل إلا سرعاتها.
\end{definition}

\begin{definition}[فوق الصوتي وتحت الصوتي]\label{def:g10:signals-and-waves:ultrasound}
تسمع أذن الإنسان \omterm{def:g10:signals-and-waves:sound}{الصوت} بين نحو \qty{20}{Hz} ونحو \qty{20}{kHz}
(ويهبط السقف مع التقدم في السنّ). \omterm{def:g10:signals-and-waves:sound}{والصوت} فوق \qty{20}{kHz} \omterm{def:g10:signals-and-waves:sound}{صوت}
\emph{فوق صوتي}\index{فوق صوتي} --- فالخفافيش تصطاد عند \qty{50}{kHz}،
والماسحات الطبية تصوّر الجسم عند بضعة ميغاهرتز. أما \omterm{def:g10:signals-and-waves:sound}{الصوت} دون
\qty{20}{Hz} فهو \emph{تحت صوتي}\index{تحت صوتي}، تحسّه الفِيَلة
ومقاييس الزلازل.
\end{definition}

\begin{proposition}[الإشارات الكهرمغناطيسية]\label{prop:g10:signals-and-waves:em}
الراديو والواي-فاي والضوء إشارات \emph{كهرمغناطيسية}\index{إشارة
كهرمغناطيسية}: أسرة واحدة، تنتقل كلها
عبر الخلاء --- لا تحتاج وسطًا --- وكلها بسرعة واحدة، هي سرعة
الضوء البالغة \cref{ch:g10:refraction}،
\[
c = \qty{3.00e8}{m/s} \quad \text{(وهي نفسها تقريبًا في الهواء).}
\]
\end{proposition}
\admitted

\begin{remark}\label{rem:g10:signals-and-waves:honest}
وهذه هنا واقعة تجريبية --- فضوء الشمس يعبر $150$ مليون
كيلومتر من الفضاء الخالي؛ ويُشتقّ اشتقاقًا أمينًا في مجلد السنة الثانية
من قوانين الكهرباء والمغناطيسية. واحفظ التقابل: \omterm{def:g10:signals-and-waves:sound}{فالصوت}
اهتزاز \emph{للوسط} ويموت بدونه؛ أما \omterm{prop:g10:signals-and-waves:em}{الإشارة الكهرمغناطيسية}
فتحمل مجالها الخاص عبر الخلاء، أسرع بمليون مرة.
\end{remark}

\begin{omfigure}
\begin{tikzpicture}[x=1.3cm, font=\small]
  \fill[omDef!15] (1.3,-0.14) rectangle (4.3,0.14);
  \draw[-Stealth] (-0.3,0) -- (7.7,0);
  \foreach \x/\l in {0/{$1$}, 1/{$10$}, 2/{$10^2$}, 3/{$10^3$},
      4/{$10^4$}, 5/{$10^5$}, 6/{$10^6$}, 7/{$10^7$}}
    \draw (\x,0.07) -- (\x,-0.07) node[below] {\l};
  \node[below] at (7.55,-0.07) {\unit{Hz}};
  \node[gray] at (0.4,-0.75) {تحت صوتي};
  \node[omDef] at (2.8,-0.75) {مسموع};
  \node[gray] at (5.8,-0.75) {فوق صوتي};
  \draw[omThm] (1.7,0.14) -- (1.7,0.4) node[above] {الشبكة};
  \draw[omThm] (2.64,0.14) -- (2.64,0.95) node[above] {نغمة لا};
  \draw[omThm] (4.7,0.14) -- (4.7,0.4) node[above] {خفاش};
  \draw[omThm] (6.7,0.14) -- (6.7,0.95) node[above] {تصوير طبي};
\end{tikzpicture}

{\small تواترات \omterm{def:g10:signals-and-waves:sound}{الصوت} على سُلّم لوغاريتمي (كل خطوة $\times 10$).
تسمع الأذن من \qty{20}{Hz} إلى \qty{20}{kHz}؛ وما دون ذلك \omterm{def:g10:signals-and-waves:ultrasound}{تحت صوتي}،
وما فوقه \omterm{def:g10:signals-and-waves:ultrasound}{فوق صوتي}.}
\end{omfigure}

\section{الأصداء: مسافات من تأخيرات}

\omterm{def:g10:signals-and-waves:signal}{الإشارة} المنتقلة بسرعة $v$ تقطع $d = v\,t$ في زمن $t$: فكل تأخير
مسافة متنكّرة. والأفضل من ذلك أن ترسل نبضة قصيرة نحو
عائق وتوقّت \emph{صداها}\index{صدى} --- فالذهاب والإياب يقطع $2d$.

\begin{method}[القياس بالصدى]\label{met:g10:signals-and-waves:echo}
\begin{enumerate}
  \item أرسِل نبضة قصيرة؛ وقِس زمن الذهاب والإياب $t$ \omterm{rem:g10:signals-and-waves:honest}{للصدى}.
  \item خذ سرعة $v$ \omterm{def:g10:signals-and-waves:signal}{الإشارة} \emph{في الوسط المقطوع}.
  \item قطعت النبضة $2d = v\,t$، ومنه $d = \dfrac{v\,t}{2}$.
\end{enumerate}
ونبضات \omterm{def:g10:signals-and-waves:sound}{الصوت} في الماء تصنع \emph{السونار}\index{سونار}
($v \approx \qty{1500}{m/s}$)؛ ونبضات الراديو في الهواء تصنع
\emph{الرادار}\index{رادار} ($v = c$)؛ ونبضات فوق \omterm{def:g10:signals-and-waves:sound}{الصوت} في الجسم تصنع
الماسحة الطبية ($v \approx \qty{1540}{m/s}$ في الأنسجة الرخوة).
\end{method}

\begin{example}[السونار]\label{ex:g10:signals-and-waves:sonar}
يرسل سونار سفينة نبضة فيعود \omterm{rem:g10:signals-and-waves:honest}{صدى} القاع بعد
$t = \qty{0.80}{s}$: فالعمق
$d = \frac{1500 \times 0.80}{2} = \qty{600}{m}$. ونسيان العامل
$2$ يضاعف المحيط.
\end{example}

\begin{omfigure}
\begin{tikzpicture}[font=\small]
  \fill[blue!8] (-3.5,0) -- (3.5,0) -- (3.5,-2.95) -- (2.6,-3.1)
    -- (1.4,-2.9) -- (0.2,-3.05) -- (-1.2,-2.85) -- (-2.3,-3.1)
    -- (-3.5,-2.95) -- cycle;
  \draw[thick, blue!60] (-3.5,0) -- (3.5,0);
  \draw[thick] (-3.5,-2.95) -- (-2.3,-3.1) -- (-1.2,-2.85)
    -- (0.2,-3.05) -- (1.4,-2.9) -- (2.6,-3.1) -- (3.5,-2.95);
  \fill[pattern=north east lines] (-3.5,-2.95) -- (-2.3,-3.1)
    -- (-1.2,-2.85) -- (0.2,-3.05) -- (1.4,-2.9) -- (2.6,-3.1)
    -- (3.5,-2.95) -- (3.5,-3.4) -- (-3.5,-3.4) -- cycle;
  \fill[gray!70] (-0.9,0.3) -- (0.9,0.3) -- (0.55,-0.25)
    -- (-0.55,-0.25) -- cycle;
  \fill[gray!50] (-0.25,0.3) rectangle (0.25,0.62);
  \draw[-Stealth, omDef, thick] (-0.35,-0.35) -- (-0.35,-2.75)
    node[midway, left] {إرسال، $t/2$};
  \draw[-Stealth, omThm, thick] (0.35,-2.75) -- (0.35,-0.35)
    node[midway, right] {صدى، $t/2$};
  \draw[Stealth-Stealth] (2.9,-0.05) -- (2.9,-2.9)
    node[midway, right] {$d = \dfrac{v\,t}{2}$};
\end{tikzpicture}

{\small نبضة سونار تنزل وتعود: فالتأخير المقيس $t$ يقطع
$2d$، ومن ثم يكون العمق $d = v\,t/2$ بأخذ $v \approx \qty{1500}{m/s}$ في
ماء البحر.}
\end{omfigure}

\begin{example}[الرادار]\label{ex:g10:signals-and-waves:radar}
يتلقى رادار مطار \omterm{rem:g10:signals-and-waves:honest}{صدى} طائرة بعد $\qty{0.30}{ms}$ من
النبضة:
$d = \frac{\qty{3.00e8}{m/s} \times \qty{3.0e-4}{s}}{2} = \qty{45}{km}$.
ولأن $c$ كبيرة جدًّا، يعيش \omterm{met:g10:signals-and-waves:echo}{الرادار} على ساعات بالميكروثانية: فتوقيت الضوء
عمل دقيق.
\end{example}

\begin{example}[القلب إشارةً]\label{ex:g10:signals-and-waves:ecg}
تبعث كل نبضة قلب نبضة توتر صغيرة عبر الصدر؛ وتسجّلها مساري
\emph{مخطط كهربية القلب}\index{مخطط كهربية القلب} (ECG)، وهو \omterm{def:g10:signals-and-waves:signal}{إشارة} شبه \omterm{def:g10:signals-and-waves:period}{دورية}.
فإذا كانت الأشواك العالية على بعد $T = \qty{0.75}{s}$ بعضها من بعض، فالقلب
ينبض \omterm{def:g10:signals-and-waves:frequency}{بتواتر} $f = 1/0.75 = \qty{1.33}{Hz}$، أي\ $1.33 \times 60 = 80$
نبضة في الدقيقة: وقراءة $T$ من أثرٍ هي كيف يأخذ طبيب القلب
نبضك.
\end{example}

\section{تمارين}

\begin{exercise}[$\star$]\label{exo:g10:signals-and-waves:1}
احسب \omterm{def:g10:signals-and-waves:frequency}{تواتر} \omterm{def:g10:signals-and-waves:period}{الإشارة الدورية} التي دورها
(أ) $T = \qty{20}{ms}$؛ (ب) $T = \qty{4.0}{ms}$؛
(ج) $T = \qty{50}{\micro s}$.
\end{exercise}

\begin{exercise}[$\star$]\label{exo:g10:signals-and-waves:2}
احسب \omterm{def:g10:signals-and-waves:period}{دور} (أ) شبكة \qty{50}{Hz}؛ (ب) نغمة لا عند \qty{440}{Hz}؛
(ج) \omterm{def:g10:signals-and-waves:signal}{إشارة} واي-فاي عند \qty{2.4}{GHz}.
\end{exercise}

\begin{exercise}[$\star$]\label{exo:g10:signals-and-waves:3}
صنّف ما يلي بين \omterm{def:g10:signals-and-waves:ultrasound}{تحت صوتي} ومسموع \omterm{def:g10:signals-and-waves:ultrasound}{وفوق صوتي}: \qty{12}{Hz}؛
و\qty{440}{Hz}؛ و\qty{18}{kHz}؛ و\qty{40}{kHz} (حسّاس ركن سيارة)؛
و\qty{5}{MHz} (مسبار طبي).
\end{exercise}

\begin{exercise}[$\star$]\label{exo:g10:signals-and-waves:4}
على راسم اهتزاز مضبوط على \qty{2}{ms} لكل تقسيمة و\qty{0.5}{V} لكل
تقسيمة، تمتد دورة واحدة على $5.0$ تقسيمات وترتفع قمة $4.0$
تقسيمة فوق المحور. جِد $T$ و$f$ \omterm{def:g10:signals-and-waves:amplitude}{والسعة}.
\end{exercise}

\begin{exercise}[$\star$]\label{exo:g10:signals-and-waves:5}
ترى البرق، ثم تسمع الرعد بعده \qty{6.0}{s}. فكم بعدت
الصاعقة ($v_{\text{الصوت}} = \qty{340}{m/s}$)؟ ولماذا يجوز إهمال
زمن انتقال الضوء؟
\end{exercise}

\begin{exercise}[$\star\star$]\label{exo:g10:signals-and-waves:6}
يعود \omterm{rem:g10:signals-and-waves:honest}{صدى} سونار من قاع البحر بعد \qty{0.60}{s}
($v = \qty{1500}{m/s}$ في ماء البحر). ولماذا يجب تنصيف الجداء $v\,t$؟
احسب العمق. وكم يستغرق \omterm{rem:g10:signals-and-waves:honest}{الصدى} فوق خندق عمقه
\qty{1200}{m}؟
\end{exercise}

\begin{exercise}[$\star\star$]\label{exo:g10:signals-and-waves:7}
يتلقى رادار مطار \omterm{rem:g10:signals-and-waves:honest}{صدى} طائرة بعد \qty{0.24}{ms}، ثم، بعد
\qty{1.0}{s} بالضبط، \omterm{rem:g10:signals-and-waves:honest}{صدى} ثانيًا بعد \qty{0.238}{ms}. احسب
المسافتين، ثم سرعة الطائرة. أهي مقتربة؟
\end{exercise}

\begin{exercise}[$\star\star$]\label{exo:g10:signals-and-waves:8}
يُطبع \omterm{ex:g10:signals-and-waves:ecg}{مخطط كهربية القلب} بسرعة \qty{25}{mm/s}؛ والأشواك العالية على بعد \qty{20}{mm}
بعضها من بعض. جِد \omterm{def:g10:signals-and-waves:period}{الدور} ومعدل ضربات القلب في الدقيقة.
وتسرّع القلب يعني أكثر من $100$ نبضة في الدقيقة: فدون أي تباعد بين الأشواك
يظهر على هذا الورق؟
\end{exercise}

\begin{exercise}[$\star\star$]\label{exo:g10:signals-and-waves:9}
عليك أن تعرض \omterm{def:g10:signals-and-waves:signal}{إشارة} تواترها \qty{200}{Hz} وسعتها \qty{3.0}{V} على
شاشة فيها $10$ تقسيمة أفقية و$8$ تقسيمة عمودية ($4$ فوق
الخط الأوسط). اختر قاعدة زمن تعرض نحو دورتين كاملتين، وكسبًا
يستعمل معظم الشاشة دون قصّ.
\end{exercise}

\begin{exercise}[$\star\star$]\label{exo:g10:signals-and-waves:10}
يرسل مسبار طبي نبضة فوق صوتية إلى البطن
($v = \qty{1540}{m/s}$ في الأنسجة الرخوة) فيسمع صدًى بعد
\qty{40}{\micro s} (الجدار الأمامي لعضو) وصدًى بعد \qty{60}{\micro s} (الجدار
الخلفي). جِد عمق كل جدار وسُمك العضو.
\end{exercise}

\begin{exercise}[$\star\star$]\label{exo:g10:signals-and-waves:11}
تُبثّ حفلة مباشرةً. فمن يسمع ضربة الطبل أولًا: مستمع
يبعد \qty{30}{m} عن المسرح، أم مستمع للراديو يبعد \qty{3000}{km}؟ وبكم
يسبقه؟
\end{exercise}

\begin{exercise}[$\star\star\star$]\label{exo:g10:signals-and-waves:12}
يصدر خفاش نبضات فوق صوتية تدوم \qty{3.0}{ms} ولا يقدر أن يسمع
صدًى وهو ما يزال يصدر. فدون أي مسافة تكون الفراشة
غير قابلة للكشف؟ وما تأخير \omterm{rem:g10:signals-and-waves:honest}{الصدى} لفراشة على بعد \qty{1.7}{m}؟ ولماذا
يجب أن يقصّر الخفاش نبضاته في الاقتراب الأخير؟
\end{exercise}

\begin{exercise}[$\star\star\star$]\label{exo:g10:signals-and-waves:13}
تصفّق مواجهًا جرفًا فتسمع \omterm{rem:g10:signals-and-waves:honest}{الصدى} بعد \qty{1.5}{s}. فكم يبعد
الجرف؟ ثم تمشي \qty{100}{m} مستقيمًا نحوه: فما التأخير الجديد؟ ولا
تعود الأذن تفصل التصفيقة عن \omterm{rem:g10:signals-and-waves:honest}{الصدى} دون نحو \qty{0.1}{s}: فضمن أي
مسافة يختفي \omterm{rem:g10:signals-and-waves:honest}{الصدى} في التصفيقة؟
\end{exercise}

\begin{exercise}[$\star\star\star$]\label{exo:g10:signals-and-waves:14}
عند \qty{5}{ms} لكل تقسيمة، تمتد دورة واحدة من \omterm{def:g10:signals-and-waves:signal}{إشارة} على $3.0$ تقسيمات
من شاشة فيها $10$ تقسيمة.
\begin{enumerate}
  \item جِد $T$ و$f$.
  \item تُحوَّل \omterm{def:g10:signals-and-waves:oscilloscope}{قاعدة الزمن} إلى \qty{2}{ms} لكل تقسيمة: فعلى كم
        تقسيمة تمتد الدورة الواحدة الآن؟
  \item وكم دورة \emph{كاملة} تسع الشاشة عند كل
        ضبط؟
\end{enumerate}
\end{exercise}

\begin{exercise}[$\star\star\star$]\label{exo:g10:signals-and-waves:15}
يقيس ``البِنغ'' زمن الذهاب والإياب لرزمة إنترنت إلى
خادوم يبعد \qty{1200}{km}. فما أصغر زمن ذهاب وإياب يمكن
تصوّره (والإشارات بسرعة $c$)؟ وفي \omterm{def:g10:refraction:fiber}{ليف بصري} ينتقل الضوء بسرعة نحو
\qty{2.0e8}{m/s}: أعِد الحساب. والبِنغ المقيس \qty{40}{ms}: أعطِ سببين
لتجاوزه جوابيك.
\end{exercise}

\section{مسألة: السونار والعاصفة ومخطط القلب}

\begin{problem}[{مسألة نهاية الأسبوع --- قراءة إشارات العالم: قانون
واحد، $d = v\,t$، يقيس بعد عاصفة، ويرسم خريطة قاع، ويعدّ نبض قلب، ويستخرج
كيف يعمل GPS}]
\label{pb:g10:signals-and-waves:1}
يعمل قارب صيد طوال ليلة عاصفة: برق في الأفق،
وسونار يمسح القاع، وقلب الربّان على شريط ورق الطبيب،
ومستقبِل GPS يصغي إلى التوابع. أربع أدوات وقانون
واحد. خذ $v_{\text{الصوت}} = \qty{340}{m/s}$ في الهواء، و\qty{1500}{m/s} في
ماء البحر، و$c = \qty{3.00e8}{m/s}$.

\medskip
\noindent\textbf{الجزء I --- العاصفة.}
\begin{enumerate}
  \item ومضة؛ ويتبعها الرعد بعد \qty{9.0}{s}. فكم تبعد؟
  \item احسب زمن انتقال الضوء على تلك المسافة، وبرّر
        معاملة الومضة معاملة الآنيّ.
  \item يعدّ البحّارة الثواني بين الومضة والرعد ويقسمون
        على ثلاثة ليحصلوا على الكيلومترات. برّر القاعدة.
  \item وبعد خمس دقائق تعطي ومضة \qty{6.0}{s}. فكم تبعد الآن؟
        جِد سرعة اقتراب العاصفة المتوسطة \omterm{def:g10:orders-of-magnitude:unit}{بوحدة} \unit{m/s} وبوحدة \unit{km/h}.
  \item يقرقع الرعد بدل أن يفرقع: فقناة البرق طولها
        كيلومترات، ومن ثم تقع أجزاؤها على مسافات مختلفة. فإذا امتدت
        القناة من \qty{2.0}{km} إلى \qty{5.0}{km} عن
        القارب، فكم تدوم القرقعة؟
\end{enumerate}

\medskip
\noindent\textbf{الجزء II --- \omterm{met:g10:signals-and-waves:echo}{السونار}.}
\begin{enumerate}[resume]
  \item يرسل \omterm{met:g10:signals-and-waves:echo}{السونار} نبضة؛ فيجيب القاع في \qty{0.90}{s}. فما العمق؟
  \item وفوق الرصيف القاري يقصر \omterm{rem:g10:signals-and-waves:honest}{الصدى} إلى \qty{0.20}{s}. فما العمق؟
  \item وتعيد نبضة واحدة \emph{صدَيين}، عند \qty{0.16}{s} وعند
        \qty{0.20}{s}. فسّرهما؛ وكم يسبح سرب السمك فوق القاع؟
  \item يرسل \omterm{met:g10:signals-and-waves:echo}{السونار} نبضة كل \qty{0.50}{s}. فما أكبر عمق
        يقدر أن يقيسه دون التباس، وما الذي يفسد وراءه؟
  \item وفي الهواء، أيّ مسافة كان سيعنيها التأخير نفسه \qty{0.60}{s}؟
        والحكمة: ما الذي يجب أن تعرفه قبل تحويل تأخير إلى
        مسافة؟
\end{enumerate}

\medskip
\noindent\textbf{الجزء III --- مخطط القلب.}
\begin{enumerate}[resume]
  \item أشواك مخطط قلب الربّان على بعد \qty{0.86}{s} بعضها من بعض. فما \omterm{def:g10:signals-and-waves:frequency}{التواتر}؟
        وكم نبضة في الدقيقة؟
  \item يتقدم الشريط بسرعة \qty{25}{mm/s}. فأيّ تباعد بين الأشواك قاسه
        الطبيب؟
  \item وبعد جرّ الشباك يصير التباعد \qty{15}{mm}. فما معدل النبض الجديد؟
  \item ويستريح بحّار السطح الشاب، وهو مجدّف مدرَّب، عند $50$ نبضة في
        الدقيقة. فما \omterm{def:g10:signals-and-waves:period}{الدور}؟ وما التباعد على الشريط؟
  \item وهل مخطط القلب دوري بدقة؟ فما الذي يقرؤه الطبيب إذن
        من الأثر؟
\end{enumerate}

\medskip
\noindent\textbf{الجزء IV --- العودة إلى الديار بالضوء.}
\begin{enumerate}[resume]
  \item يراسل القارب الميناء لاسلكيًّا، وهو يبعد \qty{50}{km}. فكم تستغرق
        الرسالة؟ وكم كان \omterm{def:g10:signals-and-waves:sound}{الصوت} سيستغرق؟
  \item يدور \omterm{def:g10:universal-gravitation:satellite}{تابع} GPS على نحو \qty{20200}{km} فوق الرأس. فكم تستغرق
        إشارته لتبلغ القارب؟
  \item يحوّل GPS الزمن إلى موضع. فأيّ خطأ في المسافة يسبّبه خطأ
        في الساعة قدره \qty{1.0}{\micro s}؟ وأيّ دقة توقيت يقتضيها تحديد
        بدقة \qty{3.0}{m}؟
  \item ولماذا لا يمكن وجود GPS يقوم على \omterm{def:g10:signals-and-waves:sound}{الصوت}، ولو من حيث المبدأ؟
  \item الخاتمة: أحصِ السرعات الثلاث المستعملة هذه الليلة والقانون الواحد
        وراء الأدوات الأربع، ثم اذكر عادة القياس \omterm{rem:g10:signals-and-waves:honest}{بالصدى} في جملة
        واحدة.
\end{enumerate}
\end{problem}
